\documentclass[aps,prl,twocolumn,superscriptaddress,floatfix]{revtex4-2}
\usepackage{amsmath,amssymb}
\usepackage{booktabs}
\usepackage{graphicx}
\usepackage{xcolor}
\usepackage{enumitem}

\newcommand{\tauexp}{\tau_{>1.2}}
\newcommand{\Tlock}{T_{\mathrm{lock}}}
\newcommand{\Dtwo}{D_2}
\newcommand{\kcat}{k_{\mathrm{cat}}}
\newcommand{\Kd}{K_d}
\newcommand{\Gtot}{G_{\mathrm{total}}}
\newcommand{\Rw}{R_{\mathrm{w}}}

\begin{document}

\title{Paper~5 Research Plan: Topological Buffering of Transient\\
Dynamical Complexity in Prebiotic Reaction Networks}
\date{\today}

\maketitle

%% ----------------------------------------------------------------
\section{Objective}

Determine whether reaction-network topology buffers the fragility
of high-dimensional transient dynamics under parameter drift.
Paper~4 showed that expanding the mutable parameter space from
1D to 3D accelerated neutral decay ($\tauexp = 0.55$ vs.\ $1.08$,
with $12.9$ lethals per population in V3b).  Paper~5 asks whether
specific network motifs can \emph{widen the evolutionary ridge}---the
region of parameter space that sustains transient dimensional
inflation---thereby mitigating this mutational hazard.

The central hypothesis:

\emph{Network motifs that create redundant slow-variable pathways
widen the viable region of parameter space for transient chaos,
reducing the effective mutational hazard quantified in Paper~4's V3b.}

%% ----------------------------------------------------------------
\section{Model Framework}

\subsection{Base System}

The Paper~3/4 enzyme-complex coupled Brusselator: two Brusselator
oscillator cores sharing a slow energy pool $E$, gated by the
catalytic complex $G/GE$.  Dynamic species:
$\{X_1, Y_1, X_2, Y_2, E\}$ (5 variables projected for $\Dtwo$).
14 reactions, all mass-action.  This is topology \textbf{T0}
(the baseline).

Fixed parameters carried from Paper~4: $A = 1.0$, $B = 3.0$,
$k_{\mathrm{on}} = k_{\mathrm{off}} = 10.0$
($\Kd = 1.0$), $\Gtot = 1.0$.

Integration: LSODA, $t \in [0, 20{,}000]$, 40\,000 points,
rtol $= 10^{-6}$, atol $= 10^{-9}$, max\_step $= 10.0$.

\subsection{Topology Augmentation Protocol}
\label{sec:embedding}

This is the critical methodological decision.  All topology
augmentations follow a \textbf{standard embedding protocol} to
ensure fair comparison:

\begin{enumerate}
  \item All new species couple to the existing system through the
    shared energy pool $E$, maintaining the slow--fast architecture
    that generates transients.

  \item All new reactions are mass-action, with rate constants
    drawn from the same prior distributions used in Phase~0
    sampling (see \S\ref{sec:phase0}).

  \item New species initial conditions are drawn from the same
    scale as existing species ($\mathcal{O}(1)$).

  \item The $\Dtwo$ projection is computed on all dynamic species
    (baseline 5 plus any additions).  This means augmented
    topologies have higher embedding dimension.
\end{enumerate}

The embedding protocol ensures that the ``interface'' between
the motif and the base system is identical across all topologies.
Differences in $\tauexp$ and ridge width are therefore attributable
to \emph{how} the motif is wired, not to coupling mechanism.

%% ----------------------------------------------------------------
\section{Topology Panel}
\label{sec:panel}

Eight topology classes, spanning three structural groups:

\subsection{Group A: Baseline and Controls}

\begin{description}
  \item[T0: Baseline.] The unmodified coupled Brusselator.
    5 species, 14 reactions.

  \item[T1a, T1b: Random-wiring controls.] Same number of added
    species and reactions as the largest motif in Group~B, but
    with random connectivity (drawn uniformly from the space of
    directed graphs with matching node and edge count).  Two
    independent random draws to avoid single-draw artifacts.
\end{description}

\subsection{Group B: Candidate Buffering Motifs}

Each motif adds 2--3 species and 3--5 reactions to T0, coupled
through the energy pool per the embedding protocol
(\S\ref{sec:embedding}).

\begin{description}
  \item[T2: Simple 3-cycle.]
    $A \to B \to C \to A$.  Minimal cyclic redundancy.

  \item[T3: Parallel dual-path cycle.]
    $A \to B \to D \to A$ and $A \to C \to D \to A$.
    Two routes through the same loop.

  \item[T4: Branch-and-rejoin.]
    $A \to B \to D$, $A \to C \to D$, $D \to A$.
    Branching with closure.

  \item[T5: Feed-forward with return.]
    $A \to B \to C$, $A \to C$, $C \to A$.
    Tests whether a short-circuit path stabilises or destabilises.

  \item[T6: Nested dual loop.]
    $A \to B \to C \to A$, $B \to D \to B$.
    Local redundancy nested inside a larger cycle.
\end{description}

For each topology, the following structural descriptors are
computed before any simulations:

\begin{itemize}
  \item Total species count and reaction count
  \item Cyclomatic number: $\mu = M - N + C$
    ($M$ = edges, $N$ = nodes, $C$ = connected components)
  \item Weak reversibility (yes/no): every reaction lies on
    some directed cycle
  \item Network deficiency (computed on the full augmented
    stoichiometric network, not the motif fragment alone;
    see \S\ref{sec:crnt_note})
\end{itemize}

\subsection{Note on CRNT Deficiency}
\label{sec:crnt_note}

Deficiency is computed on the complete augmented reaction network
(base $+$ motif), not the motif subgraph in isolation.  Because
the Brusselator model includes quasi-steady-state enzyme gating
(not pure mass-action throughout), classical CRNT deficiency
theorems do not directly apply.  Deficiency is used here as a
\textbf{structural descriptor to test}, not a theorem to invoke.
If deficiency correlates with ridge width, that is an empirical
finding requiring dynamical explanation.

%% ----------------------------------------------------------------
\section{Primary Metric: Ridge Width}
\label{sec:ridge_width}

The primary comparison metric is ridge width $\Rw$:

\begin{equation}
\Rw = \frac{\#\{\text{parameter draws with } \tauexp > 0\}}
           {\#\{\text{total draws}\}}
\end{equation}

This is the fraction of randomly sampled parameter space that
sustains any transient dimensional inflation.  A stricter variant
$\Rw^{(k)}$ (requiring $\tauexp > k$ for threshold $k = 1, 2, 3$)
is recorded as a secondary metric.

To maintain comparability with Paper~4, the mutable parameter set
for the main ridge-width comparison is $(\gamma, J, \kcat)$---the
same three parameters from V3b.  Rate constants of new motif
reactions are drawn randomly and then \emph{held fixed} during
each parameter draw.  This means we are measuring ``how robust
is this topology to drift in the energy-balance parameters?''
not ``how robust is it to drift in everything?''

A supplementary analysis varies all parameters (including new
reaction rates) and reports $\Rw$ alongside the parameter
dimensionality $d$.

%% ----------------------------------------------------------------
\section{Phenotype Definition}

Identical to Paper~4.  $\tauexp$: number of 2\,500-unit
non-overlapping sliding windows in $t \in [2{,}500, 20{,}000]$
(7~windows total) with $\Dtwo > 1.2$.  Same Grassberger--Procaccia
estimator, same Theiler window correction, same threshold.  No
metric substitutions.

$\Tlock$ is recorded as a secondary diagnostic but never used
as a selection criterion or comparison metric.

%% ----------------------------------------------------------------
\section{Numerical Safeguards}

Carried forward from Paper~4 without modification:

\begin{itemize}
  \item $\gamma$ bounds: $[\gamma_{\min} = 10^{-4},\;
    \gamma_{\max} = 0.1]$ with reflection.
  \item Solver failure $\Rightarrow \tauexp = 0$.
  \item NaN/Inf guard $\Rightarrow \tauexp = 0$.
  \item Pathological concentrations ($> 10^6$) $\Rightarrow
    \tauexp = 0$.
  \item Stiff-regime step cap: max\_step $= 10.0$.
  \item Wall-clock timeout: 120\,s per evaluation (increased
    from 60\,s to accommodate larger networks).
  \item Maximum $\tauexp = 7$.
\end{itemize}

%% ----------------------------------------------------------------
\section{Predefined Success Criteria}
\label{sec:success}

\begin{enumerate}
  \item \textbf{Positive (topology matters):}
    At least one Group~B motif has $\Rw > 2\times$ the mean
    $\Rw$ of the random-wiring controls (T1a, T1b), and the
    difference survives a permutation test at $p < 0.05$
    (10\,000 permutations of topology labels within the
    parameter-draw pool).

  \item \textbf{Strong positive (design principle):}
    In addition, a computable structural descriptor (cyclomatic
    number, deficiency, or weak reversibility) correlates with
    $\Rw$ across the panel (Spearman $|\rho| > 0.7$, $p < 0.05$).

  \item \textbf{Negative:}
    No Group~B motif exceeds the random-wiring controls by
    $> 50\%$, or no motif exceeds T0 baseline.  Reported as:
    \emph{topology does not buffer transient fragility in this
    model class.}

  \item \textbf{Dimensionality artefact:}
    Group~B motifs outperform T0 but do not outperform the
    random-wiring controls.  Interpreted as: the effect is
    driven by added species/reactions, not specific wiring.
\end{enumerate}

%% ----------------------------------------------------------------
\section{Phased Experimental Plan}

%% ---- PHASE 0 ----
\subsection{Phase~0: Topology Construction \& Structural Census}
\label{sec:phase0}

\textbf{Goal.}  Build all 8 topology classes, implement the
embedding protocol, and compute structural descriptors.  No
simulations.

\begin{enumerate}[label=0.\arabic*]
  \item Implement each topology as a variant of the Paper~4
    model file (\texttt{pilot5b\_enzyme\_complex.py}), with
    the additional species and reactions appended.

  \item For each topology, compute and tabulate: species count,
    reaction count, cyclomatic number, weak reversibility,
    deficiency.

  \item Generate 2 independent random-wiring instantiations
    (T1a, T1b) matched for species and reaction count to the
    largest Group~B motif.

  \item \textbf{Validation:}  Run 10 integration sanity checks
    per topology at the Paper~4 baseline parameters
    ($\gamma_0 = 0.00223$, $J = 4.388$, $\kcat = 0.278$)
    to confirm numerical stability.
\end{enumerate}

\textbf{Estimated cost:} $8 \times 10 = 80$ integrations.  Local.

\textbf{Decision point:}  If any topology cannot produce stable
integrations at baseline parameters, debug or remove it before
proceeding.

%% ---- PHASE 1 ----
\subsection{Phase~1: Ridge-Width Screen (All Topologies)}
\label{sec:phase1}

\textbf{Goal.}  Measure $\Rw$ for all 8 topologies via random
parameter sampling.  This is the cheap screen that determines
which topologies are promoted to full testing.

\begin{enumerate}[label=1.\arabic*]
  \item For each of the 8 topologies, draw
    $N_{\mathrm{sample}} = 300$ parameter vectors:
    \begin{align}
      \gamma &\in [0.0005,\; 0.01], \quad \text{log-uniform} \\
      J      &\in [2,\; 12], \quad \text{uniform} \\
      \kcat  &\in [0.05,\; 0.8], \quad \text{uniform}
    \end{align}
    Rate constants for motif reactions are drawn once per
    topology instantiation and held fixed across all 300 draws.
    To average over motif-rate variation, repeat with 3
    independent motif-rate draws per topology (900 total
    parameter evaluations per topology).

  \item Each parameter draw is evaluated at 2~seeds
    (seeds 42, 179, matching Paper~4).

  \item \textbf{Metrics per topology:}
    \begin{itemize}
      \item $\Rw$: fraction of draws with $\tauexp > 0$
      \item $\Rw^{(2)}$: fraction with $\tauexp > 2$
      \item Mean $\tauexp$ (across all draws, including zeros)
      \item Mean $\tauexp$ conditional on $\tauexp > 0$
    \end{itemize}

  \item \textbf{Deliverables:}
    \begin{itemize}
      \item Bar chart: $\Rw$ by topology, annotated with
        structural descriptors (Fig.~1 of paper)
      \item Scatter plots: $\Rw$ vs.\ cyclomatic number,
        $\Rw$ vs.\ deficiency
    \end{itemize}
\end{enumerate}

\textbf{Estimated cost:}
$8 \text{ topologies} \times 900 \text{ draws} \times 2 \text{ seeds}
= 14{,}400$ integrations.

At $\sim$65\,s per integration on t3.micro:
$14{,}400 \times 65 = 936{,}000\text{\,s} \approx 260\text{\,h}$.

With 10 nodes: $\sim$26\,h elapsed.

\textbf{Node assignment:}  Each node runs 1--2 topologies
(see \S\ref{sec:aws}).

\textbf{Decision point:}  If no Group~B motif exceeds the
random controls by $> 50\%$, the result is negative
(\S\ref{sec:success}, criterion~3).  Report and stop.
If Group~B motifs exceed T0 but not random controls, the result
is a dimensionality artefact (criterion~4).  Report and stop.

%% ---- PHASE 2 ----
\subsection{Phase~2: Drift-Buffering Test (Promoted Topologies)}
\label{sec:phase2}

\textbf{Goal.}  For the promoted topologies, measure how far
kinetic parameters can drift before transient complexity
collapses.

\textbf{Promotion rule:}  From Phase~1, promote 4 topologies:
\begin{itemize}
  \item Widest-ridge Group~B motif
  \item Second-widest Group~B motif
  \item Narrowest-ridge Group~B motif
  \item Best random-wiring control (T1a or T1b)
\end{itemize}

\begin{enumerate}[label=2.\arabic*]
  \item For each promoted topology, identify the parameter
    region with highest $\tauexp$ from Phase~1 data.

  \item Starting from the best parameter set, apply controlled
    drift: perturb all three parameters $(\gamma, J, \kcat)$
    simultaneously by increasing magnitudes.  Perturbation
    protocol: multiplicative noise in each parameter,
    $\theta' = \theta \cdot 10^{\mathcal{N}(0, \sigma_d^2)}$
    for $\sigma_d \in \{0.05, 0.1, 0.15, 0.2, 0.3, 0.5\}$.

  \item At each drift magnitude $\sigma_d$, run 100 perturbed
    draws $\times$ 2 seeds.

  \item \textbf{Metric:}  Survival probability
    $S(\sigma_d) = \Pr[\tauexp > 0 \mid \sigma_d]$.
    Plot $S(\sigma_d)$ vs.\ $\sigma_d$ for all 4 topologies
    on the same axes (Fig.~3 of paper).

  \item \textbf{Also measure:}  Mean $\tauexp$ conditional
    on survival; critical drift $\sigma_d^*$ at which
    $S < 0.5$.
\end{enumerate}

\textbf{Estimated cost:}
$4 \text{ topologies} \times 6 \text{ drift levels} \times
100 \text{ draws} \times 2 \text{ seeds} = 4{,}800$ integrations.
$\sim$87\,h serial, $\sim$9\,h on 10 nodes.

%% ---- PHASE 3 ----
\subsection{Phase~3: Evolutionary Robustness (Promoted Topologies)}
\label{sec:phase3}

\textbf{Goal.}  Run Paper~4-style V3b multi-parameter evolution
on the widest-ridge and narrowest-ridge promoted topologies, plus
the random-wiring control.  Show that wider-ridge topologies
exhibit less neutral decay.

\begin{enumerate}[label=3.\arabic*]
  \item \textbf{Setup:}  Identical to Paper~4 V3b.
    \begin{itemize}
      \item Population: $N = 20$.
      \item Mutable parameters: $(\gamma, J, \kcat)$.
      \item Mutation: same operators and bounds as Paper~4.
      \item Selection: tournament, $k = 3$.
      \item Generations: $G_{\max} = 40$.
      \item Seeds per evaluation: 2.
      \item Replicates: 5~selection $+$ 5~neutral per topology.
    \end{itemize}

  \item Run on 3 topologies: widest-ridge, narrowest-ridge,
    random control.

  \item \textbf{Metrics:}
    \begin{itemize}
      \item Gen-40 mean $\tauexp$ (selection vs.\ neutral)
      \item Gen-40 lethals per population (neutral condition)
      \item Cohen's $d$ between selection and neutral
    \end{itemize}

  \item \textbf{Key prediction:}  The widest-ridge topology
    will show less neutral decay (higher neutral $\tauexp$
    at gen~40, fewer lethals) than the narrowest-ridge topology.
    This directly demonstrates that topology buffers the V3b
    mutational hazard.
\end{enumerate}

\textbf{Estimated cost:}
$3 \text{ topologies} \times 10 \text{ replicates} \times
40 \text{ gen} \times 20 \text{ pop} \times 2 \text{ seeds}
= 48{,}000$ integrations.
At $\sim$65\,s: $\sim$867\,h serial, $\sim$87\,h on 10 nodes
($\sim$3.6 days).

\textbf{Node assignment:}  Each topology gets 3--4 nodes
(see \S\ref{sec:aws}).

%% ---- PHASE 4 ----
\subsection{Phase~4: Structural Predictor Analysis}
\label{sec:phase4}

\textbf{Goal.}  Test whether a computable topological property
predicts ridge width.  This phase is exploratory analysis of
Phase~1 data---no new simulations.

\begin{enumerate}[label=4.\arabic*]
  \item Compute Spearman rank correlation between $\Rw$ and
    each structural descriptor (cyclomatic number, deficiency,
    weak reversibility) across the 8 topologies.

  \item If any correlation achieves $|\rho| > 0.7$ with
    $p < 0.05$, report as a candidate design principle.

  \item If no single descriptor correlates, test simple
    combinations (e.g., cyclomatic number $\times$ weak
    reversibility indicator).

  \item \textbf{Framing:}  Explicitly exploratory.  8 topologies
    is a small sample; any correlation is a hypothesis for
    future work, not a proven law.
\end{enumerate}

\textbf{Estimated cost:} Zero additional simulations.

%% ---- PHASE 5 ----
\subsection{Phase~5: Structural Redundancy Test (Optional)}
\label{sec:phase5}

\textbf{Goal.}  For the widest-ridge topology, measure structural
redundancy by single-edge ablation.

\begin{enumerate}[label=5.\arabic*]
  \item For each reaction in the augmented network, remove it
    and re-measure $\Rw$ (using 100 parameter draws $\times$
    2 seeds from the Phase~1 protocol).

  \item Identify which reactions are load-bearing (removal
    collapses $\Rw$) vs.\ redundant (removal has minimal effect).

  \item This provides a cheap proxy for structural mutation
    robustness without requiring a full evolutionary run with
    graph-mutation operators.
\end{enumerate}

\textbf{Estimated cost:}  Depends on reaction count of the
winning topology.  If $\sim$20 reactions:
$20 \times 100 \times 2 = 4{,}000$ integrations
($\sim$72\,h serial, $\sim$7\,h on 10 nodes).

%% ----------------------------------------------------------------
\section{Computational Infrastructure}
\label{sec:aws}

\subsection{AWS Configuration}

Matching Paper~4's V2 deployment:

\begin{itemize}
  \item 10 EC2 spot instances, t3.micro, us-west-2.
  \item Spot bid: \$0.05/hr.  Expected cost: \$0.013/hr/node.
  \item Ubuntu 22.04 LTS + 2GB swap.
  \item Same deployment tooling: \texttt{deploy\_v2.sh},
    \texttt{node\_controller.py}, \texttt{monitor.py}.
  \item Project tag: \texttt{astro-paper5}.
  \item Auth token: per-experiment random token.
  \item Auto-terminate: 96\,h safety net.
\end{itemize}

\subsection{Node Assignment by Phase}

\begin{table}[h]
\centering
\caption{Node assignment across phases.  Phases are run
sequentially; nodes are reassigned between phases.}
\label{tab:nodes}
\begin{tabular}{llll}
\toprule
Phase & Task & Nodes & Elapsed \\
\midrule
0  & Topology construction  & Local  & 1 day \\
1  & Ridge-width screen     & 10     & $\sim$26\,h \\
   & \quad (analysis \& promotion) & Local & 1 day \\
2  & Drift-buffering test   & 10     & $\sim$9\,h \\
3  & Evolutionary robustness & 10    & $\sim$87\,h \\
4  & Predictor analysis     & Local  & 1 day \\
5  & Edge ablation (opt.)   & 10     & $\sim$7\,h \\
\midrule
   & \textbf{Total elapsed} &        & $\sim$8--10 days \\
   & \textbf{Total compute} &        & $\sim$1{,}200 core-h \\
   & \textbf{Est.\ AWS cost} &       & $\sim$\$15--20 \\
\bottomrule
\end{tabular}
\end{table}

Phase~1 node assignment (10 nodes, 8 topologies):
Nodes~0--7 each run one topology.  Nodes~8--9 run the two
topologies with the most draws (T1a and T1b random controls
get extra draws for statistical power).

Phase~3 node assignment (10 nodes, 3 topologies $\times$
10 replicates):  Nodes~0--3 run widest-ridge topology
(reps 0--3 selection, then 0--3 neutral via auto-advance).
Nodes~4--6 run narrowest-ridge.  Nodes~7--9 run random
control.  Remaining replicates (rep~4 for each condition)
are queued as second jobs and auto-advance after the first
completes.

\subsection{Software}

\begin{itemize}
  \item Base: \texttt{dimensional\_opening} library (existing,
    219-test-hardened).
  \item Paper~4 code: \texttt{phase1\_evolution\_v2.py} (reused
    for Phase~3 evolutionary runs).
  \item New code: \texttt{paper5\_topology\_library.py}
    (topology definitions and embedding protocol),
    \texttt{paper5\_phase1\_screen.py} (Phase~1 random sampling),
    \texttt{paper5\_phase2\_drift.py} (Phase~2 perturbation test).
  \item Output format: per-topology JSONL (matching Paper~4
    conventions), crash-resumable.
\end{itemize}

%% ----------------------------------------------------------------
\section{Decision Points}

\begin{enumerate}
  \item \textbf{After Phase~0:}  If any topology produces
    numerical instability at baseline parameters, debug or
    remove it.  If more than 2 topologies fail, revisit the
    embedding protocol.

  \item \textbf{After Phase~1:}  Apply promotion rule and
    success criteria.  If result is negative or dimensionality
    artefact, report and stop.  Do not proceed to Phase~3
    evolutionary runs.

  \item \textbf{After Phase~2:}  If drift-buffering curves are
    indistinguishable across promoted topologies, the Phase~1
    $\Rw$ differences do not translate to dynamical robustness.
    Investigate before committing to Phase~3 compute.

  \item \textbf{After Phase~3:}  If the widest-ridge topology
    does not show less neutral decay than the narrowest, the
    evolutionary-buffering claim is not supported.  Report
    Phase~1--2 results without the evolutionary extension.
\end{enumerate}

%% ----------------------------------------------------------------
\section{Scope Guardrails}

\begin{itemize}
  \item No graph-mutation evolutionary operators.  Topology is
    fixed within each experiment; only kinetic parameters mutate.
    Structural mutation (edge addition/removal as heritable
    mutations) is deferred to future work.

  \item No second oscillator model.  All experiments use
    Brusselator variants.  A Selkov cross-check is noted as a
    future direction but is out of scope.

  \item No more than 10 topologies in Phase~1.  If the panel
    grows beyond 10, prune before running.

  \item No environmental noise or time-varying parameters.

  \item No new dynamical diagnostics beyond $\tauexp$ and $\Dtwo$.

  \item Claims are restricted to: \emph{network topology does
    (or does not) buffer the fragility of transient dynamical
    complexity under parameter drift in this model class.}

  \item Results are framed as ``selective retention of topologies''
    (differential persistence), not ``evolution of topology''
    (heritable structural variation).  Consistent with Paper~4's
    retention framework.

  \item The CRNT/deficiency analysis is explicitly exploratory
    (\S\ref{sec:phase4}).  Do not claim CRNT theorems prove
    the transient-buffering result.

  \item Do not overinterpret as evidence for or against
    abiogenesis inevitability.
\end{itemize}

%% ----------------------------------------------------------------
\section{Null Model: Random-Wiring Controls}

The random-wiring controls (T1a, T1b) are the most important
methodological element of the paper.  They address the
reviewer-critical concern:

\emph{``Any improvement in $\tauexp$ could be due to adding
species, not to specific topology.''}

For the best-performing Group~B motif, an extended random
ensemble is generated in Phase~1: 20--50 independent random
rewirings, all preserving node count, edge count, and the
energy-pool coupling interface.  The motif's $\Rw$ is compared
to the full random ensemble (permutation test, $p < 0.05$).
If the motif sits above the random-control distribution, the
effect is topological, not dimensional.

This is the direct analogue of Papers~1--2's ``random vs.\
aligned growth'' comparison.

%% ----------------------------------------------------------------
\section{Expected Paper Structure}

\begin{enumerate}
  \item \textbf{Figure~1:}  Topology panel with structural
    descriptors (Table).

  \item \textbf{Figure~2:}  Phase~1 ridge-width screen
    ($\Rw$ by topology, annotated with structural properties).
    \emph{This is the headline result.}

  \item \textbf{Figure~3:}  Phase~2 drift-buffering curves
    ($S(\sigma_d)$ vs.\ drift magnitude, 4 topologies).

  \item \textbf{Figure~4:}  Motif vs.\ random-wiring ensemble
    (violin plot of random-control $\Rw$ with motif highlighted).

  \item \textbf{Figure~5:}  Phase~3 evolutionary robustness
    ($\tauexp$ trajectories and neutral-decay comparison,
    3 topologies).

  \item \textbf{Figure~6 (optional):}  Structural predictor
    scatter plots ($\Rw$ vs.\ cyclomatic number, deficiency).
\end{enumerate}

\end{document}
